\documentclass{article}
\usepackage[margin=3cm]{geometry}

\usepackage{graphicx}

\usepackage[table]{xcolor}
\usepackage{tikz}
\usetikzlibrary{shapes,arrows}
\definecolor{xlightgray}{HTML}{D1DBE4}
\definecolor{xgreen}{HTML}{B9C89F}
\definecolor{xred}{HTML}{F19F9E}
\definecolor{xyellow}{HTML}{F6DA71}
\definecolor{xblue}{HTML}{ACE7EF}
\definecolor{xpurple}{HTML}{C3A4E7}
\usepackage{todonotes}
\newcommand{\francesco}[1]{\todo[inline, color=xred]{\textbf{Francesco:} #1}}
\newcommand{\roberto}[1]{{\todo[inline, color=xblue]{\textbf{Roberto:} #1}}}
\newcommand{\luca}[1]{{\todo[inline, color=xgreen]{\textbf{Luca:} #1}}}
\newcommand{\yann}[1]{\todo[inline, color=xyellow]{\textbf{Yann:} #1}}
\usepackage{booktabs,tabularx,makecell,hyperref}
\definecolor{rowgray}{RGB}{246,246,246}

\usepackage{pdflscape}
\usepackage{environ}
\NewEnviron{interface}{
\vspace{0.5em}
\par
\begin{tikzpicture}
\node[rectangle,minimum width=0.9\textwidth] (m) {\begin{minipage}{0.85\textwidth}\BODY\end{minipage}};
\draw[dashed] (m.south west) rectangle (m.north east);
\end{tikzpicture}
}

\usetikzlibrary{arrows.meta,positioning,fit,calc,decorations.pathmorphing}

% Math
\usepackage{amsmath,amssymb}
\usepackage{mathtools}

% Algorithmic
\usepackage[noend]{algpseudocode}
\usepackage{amsthm}
\usepackage{cleveref}

\usepackage{float}
\floatstyle{ruled}
\newfloat{algo}{htbp}{algo}
\floatname{algo}{Algorithm}

% --- Custom commands ---
\newcommand{\chain}{\ensuremath{\mathsf{ch}}}
\newcommand{\genesis}{\ensuremath{B_{\mathrm{genesis}}}}
\newcommand{\V}{\ensuremath{\mathcal{V}}}
\newcommand{\B}[0]{\ensuremath{\mathcal{B}}}
\newcommand{\C}[0]{\ensuremath{\mathcal{C}}}
\newcommand{\T}[0]{\ensuremath{\mathcal{T}}}
\newcommand{\J}[0]{\ensuremath{\mathcal{J}}}
\newcommand{\M}[0]{\ensuremath{\mathcal{M}}}
\newcommand{\calS}[0]{\ensuremath{\mathcal{S}}}
\newcommand{\true}[0]{\ensuremath{\mathit{true}}}
\newcommand{\false}[0]{\ensuremath{\mathit{false}}}
\newcommand{\GST}[0]{\ensuremath{\mathsf{GST}}}
\newcommand{\slot}[0]{\ensuremath{\operatorname{\mathrm{slot}}}}

\newcommand{\grayy}[1]{{\color{gray}#1}}

% Algorithmic event environment
\algdef{SE}[EVENT]{Event}{EndEvent}[1]{\textbf{at}\ #1\ \algorithmicdo}{\algorithmicend\ \textbf{event}}
\algtext*{EndEvent}

\pagestyle{plain}

\renewcommand{\paragraph}[1]{\vspace{7pt}\noindent\textbf{#1}}

\newcommand{\block}{B}
\newcommand{\Block}{\mathsf{Block}}
\newcommand{\BlockID}{\mathsf{BlockID}}
\newcommand{\View}{\mathsf{View}}

\algrenewcommand\textproc{}

\usepackage{boxedminipage}

% Theorem environments
\theoremstyle{plain}
\newtheorem{theorem}{Theorem}
\newtheorem*{theorem*}{Theorem}
\newtheorem{corollary}{Corollary}
\newtheorem{lemma}{Lemma}
\newtheorem{proposition}{Proposition}

\theoremstyle{definition}
\newtheorem{definition}[theorem]{Definition}

\theoremstyle{remark}
\newtheorem{remark}[theorem]{Remark}


\title{Portable Finalization}

\begin{document}
\maketitle

\section{Model and definitions}

We consider a system of~$n$ validators of which at most~$f$ are adversarial, where $n \geq 3f + 1$.
A \emph{quorum} is any set of at least $2n/3$ validators.
Any two quorums overlap in at least $n/3 + 1 > f$ validators, so their intersection contains at least one honest validator.

\begin{definition}[Height]\label{def:height}
\emph{Height} is a counter for the finality gadget, decoupled from slots, rounds, and epochs.
The genesis height is $0$.
Height advances by exactly~$1$ via a single mechanism: prefix justification (Definition~\ref{def:justification}).
\end{definition}

\begin{definition}[Block]\label{def:block}
A \emph{block} has a \emph{parent} (another block) and contains a set of votes (Definition~\ref{def:vote}).
The \emph{genesis block}~$B_{\mathrm{genesis}}$ has no parent.
\end{definition}

\begin{definition}[Chain and divergence]\label{def:chain}
A \emph{chain} is a sequence of blocks forming a path from the genesis block, where each block's parent is its predecessor in the sequence.
Since a chain is uniquely determined by its last block, we identify chains with their last block and use the two interchangeably.
Two chains \emph{conflict} if neither is a prefix of the other.
The \emph{divergence point} of two conflicting chains~$X$ and~$Y$ is the last block they share: the block~$B$ such that $B$ is on both chains and no child of~$B$ is on both chains.
\end{definition}

We write $B \preceq B'$ when $B$ is an ancestor of~$B'$ (or $B = B'$), and $B \sim B'$ when $B$ and $B'$ are \emph{compatible}: $B \preceq B'$ or $B' \preceq B$.

\begin{definition}[Vote]\label{def:vote}
A \emph{vote} is a signed tuple $(\mathsf{validator},\; \mathsf{height},\; \mathsf{target},\; \mathsf{kind})$ where $\mathsf{validator}$ identifies the signer, $\mathsf{height}$ is a height, $\mathsf{target}$ is a block, and $\mathsf{kind} \in \{\mathsf{prefix},\, \mathsf{finalize}\}$ distinguishes the vote type.
A vote included on a chain is verified to reference a target that actually exists on that chain.
\end{definition}

\begin{definition}[State]\label{def:state}
The \emph{state after block~$B$}, written $\sigma[B]$, is a tuple $(h,\; \mathsf{targets},\; J,\; h_j,\; F,\; P)$.
A block~$B'$ is \emph{on the chain ending at~$B$} iff $B' \preceq B$.
\begin{itemize}
  \item $h$: the \emph{current height} of the finality gadget.
  \item $\mathsf{targets}[1..n]$: the \emph{target tracker}.
        For each validator~$i$, the highest block on the chain ending at~$B$ that~$i$ voted for at height~$h$
        ($\mathsf{targets}[i] = \bot$ if~$i$ has not voted).
  \item $J$: the \emph{justified block}, produced by the most recent justification event.
  \item $h_j$: the \emph{justified height}, the height at which~$J$ was justified.
  \item $F$: the \emph{finalized block}, the most recently finalized block.
  \item $P \subseteq \{1, \ldots, n\}$: the \emph{finality participation}, the set of validators who have cast a finalize vote ($\mathsf{kind} = \mathsf{finalize}$) with $\mathsf{height} = h_j$ and $\mathsf{target} \succeq J$.
\end{itemize}
The genesis state is $\sigma[B_{\mathrm{genesis}}] = (0,\; \bot^n,\; B_{\mathrm{genesis}},\; 0,\; B_{\mathrm{genesis}},\; \emptyset)$.
\end{definition}

\begin{definition}[Prefix justification]\label{def:justification}
Justification is evaluated on the state after each block.
A block~$B'$ on the chain ending at~$B$ is \emph{justifiable at height~$h$} (with respect to $\sigma[B]$) if at least $2n/3$ validators voted for~$B'$ or a descendant of~$B'$:
\[
  \bigl|\{i \,:\, B' \preceq \mathsf{targets}[i]\}\bigr| \;\geq\; 2n/3.
\]
The \emph{justified block} is the unique highest justifiable block.
\end{definition}

\begin{definition}[Finality]\label{def:finality}
The justified block~$J$ is \emph{finalized} when $|P| \geq 2n/3$.
\end{definition}

\begin{definition}[Slashing condition]\label{def:slashing}
Given two votes~$a$, $b$ from the same validator, the validator is \emph{slashable} if
\[
  b.\mathsf{kind} = \mathsf{finalize}
  \;\;\wedge\;\;
  a.\mathsf{height} = b.\mathsf{height}
  \;\;\wedge\;\;
  a.\mathsf{target} \not\succeq b.\mathsf{target}.
\]
That is: if a validator casts a finalize vote for block~$D$ at height~$h$, then every vote it cast at height~$h$ must have a target that is a descendant of~$D$ (or $D$ itself).
We refer to a slashable validator also as an \emph{equivocator}.

This is the \emph{only} slashing condition.
There is no double-target condition: validators may cast multiple prefix votes at the same height without slashing risk.
\end{definition}

\paragraph{Honest voting behavior.}
An honest validator casts two kinds of votes, with cleanly separated roles:
\begin{enumerate}
  \item \emph{Prefix vote} ($\mathsf{kind} = \mathsf{prefix}$, justification only).
        At each height~$h$, the validator casts a vote with $\mathsf{height} = h$, $\mathsf{target} = \textsf{getHead}()$, and $\mathsf{kind} = \mathsf{prefix}$.
  \item \emph{Finalize vote} ($\mathsf{kind} = \mathsf{finalize}$, finalization + portable justification).
        After block~$J$ is justified at height~$h_j$ and the validator has seen at least one additional block on the canonical chain at height $> h_j$ confirming that~$J$ is stable (not superseded), the validator casts a vote with $\mathsf{height} = h_j$, $\mathsf{target} = J$, and $\mathsf{kind} = \mathsf{finalize}$.
\end{enumerate}
The one-block delay on the finalize vote prevents premature finalization: a Byzantine proposer who withholds votes can cause justification at a deeper-than-optimal block, but honest validators wait to see whether the next block (from an honest proposer) supersedes it before committing.

The finalize vote provides \emph{portable justification}: on any chain at height~$h_j$ that contains~$J$, it contributes a fixed-target vote for~$J$ (Section~\ref{sec:portable}).

\begin{remark}[Non-self-slashability]\label{rem:nonslash}
An honest validator is never slashable.
The finalize vote at height~$h_j$ has target $= J$.
The prefix vote at height~$h_j$ (if any) has target $= \textsf{getHead}() \succeq J_s \succeq J$.
Both are $\succeq J$ (the finalize vote's target), so neither triggers the slashing condition.
\end{remark}

\paragraph{State transitions.}
The state after a block~$B$ is derived from the state after its parent.
The block's votes are processed against the parent's state, then justification and finality are evaluated, yielding $\sigma[B]$.

\emph{Processing a vote.}
When a vote~$v$ in block~$B$ is processed, the target tracker is updated if the vote is for the current height and the target is on the chain ending at~$B$ ($v.\mathsf{target} \preceq B$): $\mathsf{targets}[v.\mathsf{validator}]$ is set to $v.\mathsf{target}$ when it is higher than the current record (or the validator has not yet voted).
This applies to both prefix and finalize votes.
Independently, if the vote is a finalize vote ($v.\mathsf{kind} = \mathsf{finalize}$) with $v.\mathsf{height} = h_j$, $v.\mathsf{target} \succeq J$, and $v.\mathsf{target} \preceq B$, the validator is added to the finality participation~$P$.

\emph{Height processing.}
After processing all votes in a block, justification and finality are evaluated.
If $|P| \geq 2n/3$, the justified block is finalized: $F \gets J$.
Then the chain is walked from~$B$ toward genesis, accumulating the count of validators whose target is each block.
The first block~$B'$ where the cumulative count reaches $2n/3$ is the highest justifiable block; if found, the height advances ($h_j \gets h$, $h \gets h+1$, targets reset).
The justified block~$J$ and finality participation~$P$ are updated only if the new justified block is strictly higher than the current one ($B' \succ J$); otherwise $J$ and $P$ are unchanged.
Finality is checked before justification, so a justified block can be finalized in the same round it is superseded.

\begin{center}
\begin{boxedminipage}{\textwidth}
\begin{algorithmic}[1]

\Function{$\textsf{processBlock}$}{$s,\; B$}
  \For{each vote $v$ in $B$} $s \gets \textsf{processVote}(s,\, v,\, B)$ \EndFor
  \State \Return $\textsf{processHeight}(s,\, B)$
\EndFunction
\vspace{3pt}

\Function{$\textsf{processVote}$}{$(h,\, \mathsf{targets},\, J,\, h_j,\, F,\, P),\; v,\; B$}
  \State $i \gets v.\mathsf{validator}$
  \If{$v.\mathsf{height} = h$ and $v.\mathsf{target} \preceq B$} \Comment{Both prefix and finalize votes}
    \If{$\mathsf{targets}[i] = \bot$ or $v.\mathsf{target} \succ \mathsf{targets}[i]$}
      \State $\mathsf{targets}[i] \gets v.\mathsf{target}$
    \EndIf
  \EndIf
  \If{$v.\mathsf{kind} = \mathsf{finalize}$ and $v.\mathsf{height} = h_j$ and $J \preceq v.\mathsf{target} \preceq B$} \Comment{Finalize votes only}
    \State $P \gets P \cup \{i\}$
  \EndIf
  \State \Return $(h,\, \mathsf{targets},\, J,\, h_j,\, F,\, P)$
\EndFunction
\vspace{3pt}

\Function{$\textsf{processHeight}$}{$(h,\, \mathsf{targets},\, J,\, h_j,\, F,\, P),\; B$}
  \If{$|P| \geq 2n/3$} \Comment{Finality}
    \State $F \gets J$
  \EndIf
  \State $\mathsf{count} \gets 0$ \Comment{Justification: walk from $B$ toward genesis}
  \For{$B'$ from $B$ down to $B_{\mathrm{genesis}}$}
    \State $\mathsf{count} \gets \mathsf{count} + |\{i : \mathsf{targets}[i] = B'\}|$
    \If{$\mathsf{count} \geq 2n/3$} \Comment{Height advances}
      \If{$B' \succ J$} \Comment{Justified block advances}
        \State $J \gets B'$, \; $h_j \gets h$, \; $P \gets \emptyset$
      \EndIf
      \State $h \gets h + 1$
      \For{all $i$} $\mathsf{targets}[i] \gets \bot$ \EndFor
      \State \Return $(h,\, \mathsf{targets},\, J,\, h_j,\, F,\, P)$
    \EndIf
  \EndFor
  \State \Return $(h,\, \mathsf{targets},\, J,\, h_j,\, F,\, P)$
\EndFunction

\end{algorithmic}
\end{boxedminipage}
\end{center}


\section{Accountable safety}

\begin{lemma}[Height progression]\label{lem:heightprog}
For any chain~$Y$ to reach height~$h' > h$, chain~$Y$ must have advanced through every intermediate height.
At each height $h'' \in \{h, \ldots, h'-1\}$, prefix justification reached a quorum on chain~$Y$.
\end{lemma}

\begin{proof}
Height advances by exactly~$1$ per justification event.
Reaching~$h'$ from~$h$ requires $h' - h$ successive justification quorums.
\end{proof}

\begin{lemma}[Justified block monotonicity]\label{lem:slotmono}
On any single chain, the justified block only moves forward: if~$J$ is the justified block at some point, any future justified block~$J'$ satisfies $J \preceq J'$.
Furthermore, $F \preceq J$ at all times.
\end{lemma}

\begin{proof}
$J$ is only updated when $B' \succ J$ (the guard in $\textsf{processHeight}$).
$F$ is set to~$J$ in $\textsf{processHeight}$; by monotonicity of~$J$, $F \preceq J$ is maintained.
\end{proof}

\begin{lemma}[Any chain past a finalized height contains the finalized block]\label{lem:mainsafety}
If $C$ is finalized at height~$h$, any chain~$Y$ that has progressed past height~$h$ must contain~$C$---otherwise $\geq n/3$ validators are slashable.
\end{lemma}

\begin{proof}
By Definition~\ref{def:finality}, finality of~$C$ at height~$h$ on chain~$X$ means $|P| \geq 2n/3$: a set~$Q_F$ of at least $2n/3$ validators where each member~$i$ signed a finalize vote with $\mathsf{height} = h$ and $\mathsf{target} = D_i \succeq C$.

By Lemma~\ref{lem:heightprog}, chain~$Y$ reaching height $h' > h$ progressed through height~$h$ via justification: a quorum~$Q$ of at least $2n/3$ validators each voted at height~$h$ on chain~$Y$.
By quorum intersection, $|Q_F \cap Q| \geq n/3$.

Consider any $v \in Q_F \cap Q$.
From~$Q_F$, $v$ signed a finalize vote at height~$h$ with target $D_v \succeq C$.
From~$Q$, $v$ voted at height~$h$ on chain~$Y$ with some target~$T_v \preceq Y$.
If $T_v \not\succeq D_v$, then $v$ is slashable (Definition~\ref{def:slashing}: the finalize vote has $\mathsf{kind} = \mathsf{finalize}$, $\mathsf{height} = h$, and $T_v \not\succeq D_v$).

Since $|Q_F \cap Q| \geq n/3$ and at most~$f < n/3$ validators are slashable, at least one $v \in Q_F \cap Q$ is non-slashable, so $T_v \succeq D_v$.
This gives $T_v \succeq D_v \succeq C$.
Since $T_v \preceq Y$ (on chain~$Y$) and $C \preceq D_v \preceq T_v$, chain~$Y$ contains~$C$.
\end{proof}

\begin{lemma}[Justifications are compatible with finalized blocks]\label{lem:justcompat}
Unless $\geq n/3$ validators are slashable: if $C$ is finalized at height~$h$, then any block justified at height $\geq h$ on any chain is compatible with~$C$.
\end{lemma}

\begin{proof}
A justification at height~$h'$ advances the chain to height $h' + 1 > h' \geq h$, so the chain has progressed past height~$h$.
By Lemma~\ref{lem:mainsafety}, that chain contains~$C$; since the chain is linear, the justified block satisfies $B \sim C$.
\end{proof}

\begin{theorem}[Accountable safety]\label{thm:safety}
Unless $\geq n/3$ validators are slashable, no two conflicting blocks can be finalized.
\end{theorem}

\begin{proof}
Suppose $C$ is finalized at height~$h$ and $C'$ at height~$h'$, with $h \leq h'$ without loss of generality.
Finality of~$C'$ at height~$h'$ requires that~$C'$ was also justified at height~$h' \geq h$.
By Lemma~\ref{lem:justcompat}, $C' \sim C$.
\end{proof}


\section{Fork-choice store}

The previous section models finality as it evolves \emph{within} a chain.
A node, however, may track multiple chains simultaneously and must select one to follow.
The \emph{store} is the node-local data structure that mediates this selection.

\begin{definition}[Store]\label{def:store}
A node maintains a \emph{store} $(\sigma,\; J_s,\; F_s,\; \mathcal{T},\; C)$ where:
$\sigma$ is the \emph{state map} (Definition~\ref{def:state}), mapping each accepted block to its post-state,
$J_s$ is the \emph{store-justified block},
$F_s$ is the \emph{store-finalized block},
$\mathcal{T}$ is the \emph{block tree}---the set of all blocks the node has accepted,
and $C$ is the \emph{candidate set}---the set of all $(J,\, h)$ pairs obtained from the post-states of processed blocks.

The genesis store has $J_s = F_s = B_{\mathrm{genesis}}$, $\mathcal{T} = \{B_{\mathrm{genesis}}\}$, and $C = \emptyset$.
\end{definition}

The store is updated when a new block is received.
The node computes $\sigma[B] \gets \textsf{processBlock}(\sigma[B.\mathsf{parent}],\; B)$, then updates the store's justified and finalized blocks.

\emph{Accepting a block.}
A block~$B$ is accepted into the store only if its parent is already in~$\mathcal{T}$ and it descends from~$F_s$.
Blocks are processed in parent-first order.
The candidate pair $(\sigma[B].J,\; \sigma[B].h_j)$ is added to~$C$.

\emph{Selecting the justified block.}
$\textsf{computeJustified}$ is a deterministic function of~$C$ and~$F_s$.
It starts at~$F_s$ and walks toward descendants: among all candidates in~$C$ whose block strictly descends from the current position, it picks the one with the greatest height (breaking ties by a fixed total order on blocks), and repeats.
When no further descendant exists in~$C$, the walk terminates and the final block is returned.

\emph{Updating finalized.}
If $F_B \succ F_s$ and $F_B \preceq J_s$, then $F_s \gets F_B$.

\begin{center}
\begin{boxedminipage}{\textwidth}
\begin{algorithmic}[1]

\State \textbf{Store:} $J_s \gets B_{\mathrm{genesis}}$, \; $F_s \gets B_{\mathrm{genesis}}$, \; $\mathcal{T} \gets \{B_{\mathrm{genesis}}\}$, \; $C \gets \emptyset$, \; $\sigma[B_{\mathrm{genesis}}] \gets (0,\; \bot^n,\; B_{\mathrm{genesis}},\; 0,\; B_{\mathrm{genesis}},\; \emptyset)$
\vspace{3pt}

\Function{$\textsf{onBlock}$}{$B$}
  \State \textbf{assert} $B.\mathsf{parent} \in \mathcal{T}$ \Comment{Parent already accepted}
  \State \textbf{assert} $F_s \preceq B$ \Comment{$B$ descends from finalized}
  \State $\mathcal{T} \gets \mathcal{T} \cup \{B\}$
  \State $\sigma[B] \gets \textsf{processBlock}(\sigma[B.\mathsf{parent}],\; B)$
  \State $C \gets C \cup \{(\sigma[B].J,\; \sigma[B].h_j)\}$
  \State $J_s \gets \textsf{computeJustified}()$
  \State $\textsf{updateFinalized}(\sigma[B].F)$
\EndFunction
\vspace{3pt}

\Function{$\textsf{computeJustified}$}{}
  \State $J^* \gets F_s$
  \While{$\exists\, (J',\, h') \in C$ with $J' \succ J^*$} \Comment{Walk toward descendants}
    \State $J^* \gets \arg\max_{(J',\, h') \in C,\; J' \succ J^*} (h',\, J')$
  \EndWhile
  \State \Return $J^*$
\EndFunction
\vspace{3pt}

\Function{$\textsf{updateFinalized}$}{$F_B$}
  \If{$F_B \succ F_s$ and $F_B \preceq J_s$}
    \State $F_s \gets F_B$
  \EndIf
\EndFunction

\end{algorithmic}
\end{boxedminipage}
\end{center}

\begin{definition}[$\textsf{getHead}$]\label{def:gethead}
The fork-choice function $\textsf{getHead}$ takes the store and returns a block from the subtree rooted at~$J_s$ in~$\mathcal{T}$.
The precise selection rule is not specified here; the only property we require is that the output always descends from~$J_s$.
\end{definition}

\subsection{Store invariants and safety}

\subsubsection*{Unconditional invariants}

\begin{theorem}[$F_s \preceq J_s$ invariant]\label{thm:fleqj}
The store maintains $F_s \preceq J_s$ at all times.
This holds without any assumption on~$f$.
\end{theorem}

\begin{proof}
$\textsf{computeJustified}$ starts the walk at~$F_s$ and only moves to strict descendants, so it returns $J_s \succeq F_s$.
$\textsf{updateFinalized}$ updates $F_s$ only when $F_B \preceq J_s$, so $F_s \preceq J_s$ is maintained.
\end{proof}

\begin{theorem}[Local irreversibility of finality]\label{thm:finperm}
Once a node sets $F_s = F$, the store-finalized block descends from~$F$ at all future times.
This holds without any assumption on~$f$.
\end{theorem}

\begin{proof}
$F_s$ is only updated in $\textsf{updateFinalized}$, which requires $F_B \succ F_s$.
\end{proof}

\begin{theorem}[Fork-choice consistency]\label{thm:fcconsistency}
Once a node sets $F_s = F$, $\textsf{getHead}$ returns a block descending from~$F$ at all future times.
This holds without any assumption on~$f$.
\end{theorem}

\begin{proof}
By Theorem~\ref{thm:finperm}, $F \preceq F_s$ at all future times.
By Theorem~\ref{thm:fleqj}, $F_s \preceq J_s$.
$\textsf{getHead}$ returns a descendant of~$J_s$, so the head descends from~$F_s \succeq F$.
\end{proof}

\subsubsection*{Properties under $f < n/3$}

\begin{lemma}[Upgrade property of $\textsf{computeJustified}$]\label{lem:upgrade}
Unless $\geq n/3$ validators are slashable: if $F$ is finalized at height~$h_F$, $(F, h_F) \in C$, and $F_s \prec F$, then $\textsf{computeJustified}$ produces $J_s \succeq F$.
\end{lemma}

\begin{proof}
$F_s$ and $F$ are both finalized, so by Theorem~\ref{thm:safety}, $F_s \sim F$.
Since $F_s \prec F$, $F$ is a strict descendant of~$F_s$, and $(F, h_F) \in C$.
The walk starts at $J^* = F_s$.
At any step where $J^* \prec F$: $(F, h_F) \in C$ and $F \succ J^*$, so the walk does not terminate.
The selected candidate has height $\geq h_F$.
By Lemma~\ref{lem:justcompat}, any justification at height $\geq h_F$ is compatible with~$F$.
A strict descendant of~$J^*$ (where $J^* \prec F$) that is compatible with~$F$ is either $\prec F$ or $\succeq F$.
Since~$C$ is finite, the walk reaches $J^* \succeq F$ in finitely many steps.
Once $J^* \succeq F$, every subsequent step moves to $J' \succ J^* \succeq F$.
When the walk terminates, $J_s = J^* \succeq F$.
\end{proof}

\begin{theorem}[Local acceptance of finality updates]\label{thm:finlive}
Unless $\geq n/3$ validators are slashable: if a block~$B$ is processed by $\textsf{onBlock}$ and its post-state has $F_B \succ F_s$, then after processing, $F_s$ is updated to~$F_B$.
\end{theorem}

\begin{proof}
We need $F_B \preceq J_s$ after $\textsf{computeJustified}$ runs, so that the guard in $\textsf{updateFinalized}$ passes.
Finalizing~$F_B$ requires a prior justification of~$F_B$: at some earlier point on the chain, a block was processed whose post-state had $(F_B, h_F)$ as its justified checkpoint.
Since~$B$ is in the block tree, every ancestor was accepted, so $(F_B, h_F) \in C$.
Since $F_B \succ F_s$, Lemma~\ref{lem:upgrade} gives $J_s \succeq F_B$.
\end{proof}

\begin{theorem}[Lock-in]\label{thm:lockin}
Unless $\geq n/3$ validators are slashable: if block~$F$ is finalized at height~$h$ on any chain, and a node has processed a block whose post-state has $(F, h)$ as its justified checkpoint, then $J_s \succeq F$ at all future times.
Consequently, $F$ is always on the node's canonical chain.
\end{theorem}

\begin{proof}
Once the node processes such a block, $(F, h) \in C$ permanently.
If $F_s \succeq F$, then $J_s \succeq F_s \succeq F$ (Theorem~\ref{thm:fleqj}).
If $F_s \prec F$, Lemma~\ref{lem:upgrade} gives $J_s \succeq F$.
\end{proof}

\begin{theorem}[Order independence]\label{thm:orderindep}
Unless $\geq n/3$ validators are slashable: the store state $(J_s, F_s)$ after processing a set of blocks depends only on \emph{which} blocks have been processed, not on the order.
\end{theorem}

\begin{proof}
The candidate set~$C$ is order-independent (adding elements to a set is commutative).
$\textsf{computeJustified}$ is a deterministic function of $(C, F_s)$, so $J_s$ depends only on~$C$ and~$F_s$.
Under $f < n/3$, all finalized blocks are mutually compatible (Theorem~\ref{thm:safety}).
Let $F_{\max}$ be the maximum among all $F_B$ values.
By Theorem~\ref{thm:finlive}, processing the block carrying $F_{\max}$ yields $F_s = F_{\max}$.
No subsequent block changes~$F_s$.
Thus both~$C$ and~$F_s$ are order-independent, and so is~$J_s$.
\end{proof}

\begin{corollary}[Cross-node agreement]\label{cor:agreement}
Unless $\geq n/3$ validators are slashable: if two nodes have processed the same set of blocks, they have the same $(J_s, F_s)$.
\end{corollary}


\section{Portable justification}\label{sec:portable}

The finalize vote has a dual role: it finalizes the justified block on the chain where the original prefix justification fired, and it provides a \emph{portable fixed-target justification} on any other chain.

\paragraph{Mechanism.}
After block~$J$ is prefix-justified at height~$h_j$ on chain~$X$ (and confirmed stable), honest validators send a finalize vote: $(\mathsf{validator},\; h_j,\; J,\; \mathsf{finalize})$.
This vote has $\mathsf{height} = h_j$, $\mathsf{target} = J$, and $\mathsf{kind} = \mathsf{finalize}$.
Consider how it is processed on different chains:

\emph{Original chain~$X$ (at height $> h_j$).}
The target part does not register ($v.\mathsf{height} = h_j \neq h$).
The finalize part registers ($v.\mathsf{kind} = \mathsf{finalize}$, $v.\mathsf{height} = h_j = h_j$, and $J \preceq v.\mathsf{target} = J \preceq B$), adding the validator to~$P$.
When $|P| \geq 2n/3$, $J$ is finalized.

\emph{Another chain~$Y$ at height~$h_j$, containing~$J$.}
The target part registers: $v.\mathsf{height} = h_j = h$ and $v.\mathsf{target} = J \preceq B$ (since $J$ is on chain~$Y$).
If $\geq 2n/3$ finalize votes are included, their targets all equal~$J$, so the prefix count at~$J$ reaches $\geq 2n/3$.
Prefix justification fires for~$J$ (since all votes target~$J$ specifically, $J$ is the highest prefix-justified block).
Height advances on chain~$Y$.

\emph{Chain~$Z$ not containing~$J$.}
Both checks fail: $v.\mathsf{target} = J \not\preceq B$ (since $J$ is not on chain~$Z$).
The finalize vote has no effect.

\paragraph{Why this provides re-org resistance.}
The slashing condition (Definition~\ref{def:slashing}) locks validators who cast finalize votes: any vote at height~$h_j$ must have a target $\succeq J$.
A validator who finalize-voted for~$J$ cannot vote for blocks incompatible with~$J$ at height~$h_j$ without being slashable.
This means:
\begin{itemize}
  \item Under $f < n/3$, any chain that advances past height~$h_j$ must contain~$J$ (Lemma~\ref{lem:mainsafety}).
  \item The finalize votes make~$J$ justified on any chain that reaches height~$h_j$ and contains~$J$, ensuring the fork-choice anchors on~$J$ everywhere.
  \item A competing chain that diverges before~$J$ cannot include the finalize votes (the on-chain verification $v.\mathsf{target} \preceq B$ fails), so it cannot benefit from them.
\end{itemize}

\begin{remark}[Contrast with pure timeouts]
In Simplex, a timeout vote carries no target ($\mathsf{target} = \bot$), so there is no block to ``replay'' on a competing chain.
If a height was advanced via timeout, other chains have no certificate to import.
Portable finalization solves this: the finalize vote carries an explicit target ($J$), making it replayable on any chain containing~$J$.
\end{remark}

\begin{remark}[Finalization on competing chains]
On chain~$Y$, after~$J$ is justified (from the finalize votes' target part in one block), those votes' finalize parts register in subsequent blocks (once the chain's $h_j$ equals~$h_j$ and $J \preceq v.\mathsf{target} \preceq B$), growing~$P$ toward finalization.
This is consistent with safety: $J$ is on both chains.
\end{remark}


\section{Inactivity leak}

We assume \emph{synchrony}: all messages between honest validators are delivered within a known bound, at least one honest validator proposes a block in each round, and $\textsf{getHead}$ is a deterministic function of the store.
A \emph{round} is one block-processing step at a given height.
The \emph{canonical chain} is the chain selected by $\textsf{getHead}$.
Under synchrony, honest validators have received the same set of blocks and therefore share the same store state (Corollary~\ref{cor:agreement}) and canonical chain.

\subsection{Penalty units}

During periods of non-finality, the protocol assesses \emph{penalty units} per validator per round.

\begin{definition}[Penalty units]\label{def:penalty}
Let $r$ denote the round index within the current height ($r = 1$ for the first round after the height last advanced).
Penalties are assessed in two layers depending on whether height advances.

\medskip\noindent\textbf{Layer 1 (stall --- height does not advance):}
$\mathsf{penalty}(i) = 0$ if $\mathsf{targets}[i] \neq \bot$ (the validator voted at this height), else $1$.

\medskip\noindent\textbf{Layer 2 (advance --- height advances):}
two independent checks, each contributing $0$ or $1$.
\begin{enumerate}
  \item \emph{Participation check}: $1$ unit unless $\mathsf{targets}[i] \neq \bot$.
  \item \emph{Finalize check} (only at $h = h_j + 1$ with $|P| < 2n/3$): $1$ unit unless $i \in P$.
\end{enumerate}
Thus $\mathsf{penalty}(i) \in \{0, 1, 2\}$ in an advance round.
\end{definition}

\subsection{Leak fairness}

\begin{theorem}[Leak fairness]\label{thm:fairness}
Under $f < n/3$ and synchrony, every honest validator has $\mathsf{penalty}(i) = 0$ on the canonical chain in every round.
\end{theorem}

\begin{proof}
Let~$i$ be an honest validator and~$X$ the canonical chain.
Under synchrony, all honest validators vote for $\textsf{getHead}()$ at the current height.
Since $\geq 2n/3$ honest validators all vote for the same block, the prefix count at that block reaches $\geq 2n/3$.
Justification fires in one round (no stalls occur), so Layer~1 penalties never apply.

For Layer~2: validator~$i$ voted ($\mathsf{targets}[i] \neq \bot$), so the participation check is satisfied.
For the finalize check: after one block of confirmation, $i$ sends a finalize vote with $\mathsf{finalize\_height} = h_j$ and $\mathsf{finalize\_target} = J$, so $i \in P$.
Under synchrony, the finalize vote is included in the next honest proposer's block, and finality fires (finality is checked before justification in $\textsf{processHeight}$, so~$P$ is evaluated before the checkpoint can advance).

Crucially, there is no ``lock'' that prevents honest validators from voting on the canonical chain.
The slashing condition requires only that votes at height~$h_j$ are compatible ($\succeq$) with the finalize target~$J$.
On the canonical chain, $\textsf{getHead}() \succeq J_s \succeq J$, so every honest vote is compatible.
\end{proof}

\subsection{Leak tightness}

\begin{theorem}[Amortized $n/3$ penalty units per height]\label{thm:tightness}
Under $f < n/3$: during any period of non-finality, the total penalty units is at least $n/3$ per height on average.
\end{theorem}

\begin{proof}
Every height is either a stall (no advance) or an advance.

\emph{Stall height (Layer~1 only).}
Let $m = |\{i : \mathsf{targets}[i] \neq \bot\}|$ be the participation count.
Since justification didn't fire, $m < 2n/3$ (otherwise the prefix count at the genesis block alone would be $\geq m \geq 2n/3$).
The penalty is $n - m > n/3$.

\emph{Advance height where the checkpoint updates (Layer~2).}
The participation penalty may be zero.
However, the checkpoint update means finalization is now pending.
At the next advance height $h = h_j + 1$ with $|P| < 2n/3$, the finalize check contributes $n - |P| > n/3$.

\emph{Advance height without checkpoint update (Layer~2).}
The justified block $B'$ has $B' \preceq J$ (the checkpoint didn't advance).
The prefix count reached $2n/3$ at $B'$, but the participation check counts validators who voted at all (not just for the canonical target).
However, $B' \preceq J$ means the honest validators' votes (for head $\succeq J \succeq B'$) all contributed to the prefix count at $B'$ via prefix counting.
So $\geq 2n/3$ validators participated, and the participation penalty is $\leq n/3$.
We attribute the $\geq n/3$ non-finalizing validators from the next finalize-check round (as above).

In every case, each height contributes $\geq n/3$ penalty units.
\end{proof}


\appendix
\section{Differences from Prefix-IC}\label{app:diffs-prefixic}

This section summarizes the changes from Prefix-IC to Portable Finalization.

\subsection*{Vote structure}

\begin{center}
\begin{tabularx}{\textwidth}{lXX}
\toprule
& \textbf{Prefix-IC} & \textbf{Portable Finalization} \\
\midrule
Fields & $(\mathsf{val},\, \mathsf{height},\, \mathsf{target},\, \mathsf{fin\_h},\, \mathsf{fin\_target})$ & $(\mathsf{val},\, \mathsf{height},\, \mathsf{target},\, \mathsf{kind})$ \\
Finalize commitment & Via $\mathsf{fin\_h}$/$\mathsf{fin\_target}$ fields on any vote & Via $\mathsf{kind} = \mathsf{finalize}$ (height and target are the commitment) \\
Vote kinds & One (finalize fields optional) & Two: $\mathsf{prefix}$ and $\mathsf{finalize}$ \\
\bottomrule
\end{tabularx}
\end{center}

The finalize fields ($\mathsf{finalize\_height}$, $\mathsf{finalize\_target}$) are replaced by a $\mathsf{kind}$ field.
For a finalize vote, the vote's own $\mathsf{height}$ and $\mathsf{target}$ serve as the finalize height and finalize target, eliminating redundancy.

\subsection*{Slashing condition}

\begin{center}
\begin{tabularx}{\textwidth}{lXX}
\toprule
& \textbf{Prefix-IC} & \textbf{Portable Finalization} \\
\midrule
Condition & $d_1.\mathsf{target} \neq d_2.\mathsf{fin\_target}$ (equality) & $a.\mathsf{target} \not\succeq b.\mathsf{target}$ where $b.\mathsf{kind} = \mathsf{finalize}$ (compatibility) \\
Effect & Vote target must exactly match finalize target & Vote target must descend from finalize vote's target \\
\bottomrule
\end{tabularx}
\end{center}

Two changes: (1)~the relaxation from equality to compatibility, enabling validators to vote for head $\succeq J$ while finalizing~$J$; (2)~the condition references $\mathsf{kind}$ instead of separate finalize fields.

\subsection*{Honest voting behavior}

\begin{center}
\begin{tabularx}{\textwidth}{lXX}
\toprule
& \textbf{Prefix-IC} & \textbf{Portable Finalization} \\
\midrule
Prefix vote & target $=$ head, finalize fields optional & $\mathsf{kind} = \mathsf{prefix}$, target $=$ head \\
Finalize vote & --- (finalize fields on prefix vote) & $\mathsf{kind} = \mathsf{finalize}$, target $= J$, height $= h_j$ \\
Timing & Immediate & One-block delay (timer) \\
\bottomrule
\end{tabularx}
\end{center}

Prefix votes and finalize votes have cleanly separated roles and distinct vote structures.
The finalize vote's dual role (finalization on the original chain, portable justification on other chains) emerges from the independent processing of the target and finality-participation checks in $\textsf{processVote}$.
The one-block delay prevents a Byzantine proposer from causing premature finalization of a manipulated justified block.

\subsection*{State machine and store}

Both are unchanged from Prefix-IC.
The $\textsf{processVote}$, $\textsf{processHeight}$, $\textsf{computeJustified}$, and $\textsf{updateFinalized}$ functions are identical.
All store invariants and safety proofs carry over with the adapted slashing condition.

\subsection*{New: portable justification (Section~\ref{sec:portable})}

The finalize vote provides fixed-target justification for~$J$ on any chain at height~$h_j$ that contains~$J$.
This is the main new feature.
It ensures the fork-choice anchors on~$J$ everywhere, preventing adversarial re-orgs during sleepiness.

\subsection*{No timeouts}

Unlike Simplex and Rooted Timeouts, this protocol has no timeout or notarization mechanism.
Heights advance exclusively via prefix justification.
The compatibility-based slashing condition eliminates the ``lock'' problem: honest validators are never prevented from voting on the canonical chain (Theorem~\ref{thm:fairness}), so there is no need for a timeout mechanism to bypass locks.

\end{document}
